\chapter{Cushing's Syndrome}

\section*{Definition and Overview}
Cushing's syndrome is a condition caused by prolonged exposure to high levels of the hormone cortisol. It produces a characteristic pattern of weight gain, particularly around the face, neck, and abdomen, along with thinning skin, easy bruising, high blood pressure, and high blood sugar. The most common cause is long-term use of corticosteroid medicines (iatrogenic Cushing's), while other causes include tumours of the adrenal or pituitary glands that produce excess cortisol. Cushing's syndrome is serious but usually treatable, and outcomes are good when the cause is found and corrected.

\section*{Epidemiology}
Cushing's syndrome is uncommon, affecting about one to two people per 100,000 per year, and is more common in women than men, with a peak in adults aged 20--50. The use of steroid medicines for conditions such as asthma, arthritis, and inflammatory disease is the most frequent cause overall. Tumour-related Cushing's is rare: the most common is Cushing's disease, caused by a benign pituitary tumour, which accounts for about 70\% of non-iatrogenic cases. Because the features develop gradually, the diagnosis is often delayed for years.

\section*{Aetiology and Pathophysiology}
Excess cortisol can come from external medicines or from the body's own production.

\begin{itemize}
    \item \textbf{Exogenous (medicines):} long-term corticosteroid treatment, including tablets, injections, and some creams and inhalers at high doses
    \item \textbf{Cushing's disease:} a benign pituitary tumour overproduces ACTH, which stimulates the adrenal glands to make cortisol
    \item \textbf{Adrenal tumours:} benign or, rarely, malignant tumours of the adrenal gland produce cortisol directly
    \item \textbf{Ectopic ACTH:} tumours elsewhere in the body, such as in the lung, produce ACTH
    \item \textbf{Effects of cortisol:} it raises blood sugar, redistributes fat, breaks down protein, and suppresses the immune system
    \item \textbf{Feedback loss:} normal cortisol regulation is overwhelmed by the excess
\end{itemize}

\section*{Clinical Features}
The features develop gradually and are often mistaken for ordinary weight gain or ageing.

\begin{itemize}
    \item Weight gain, especially around the face (moon face), neck (buffalo hump), and abdomen
    \item Thinning skin and easy bruising
    \item Purple stretch marks on the abdomen, thighs, and breasts
    \item Muscle weakness, particularly of the legs
    \item High blood pressure and high blood sugar or diabetes
    \item Round, red face
    \item Mood changes, anxiety, depression, and irritability
    \item Irregular periods and reduced libido
    \item Slowed growth in children
    \item Increased risk of infections and poor wound healing
\end{itemize}

\section*{Differential Diagnosis}
Many common conditions produce similar features.

\begin{itemize}
    \item Obesity and metabolic syndrome
    \item Polycystic ovary syndrome, with weight gain and irregular periods
    \item Type 2 diabetes and high blood pressure
    \item Depression and anxiety
    \item Hypothyroidism
    \item Alcoholism, which can cause a similar appearance
    \item Simple weight gain from lifestyle
\end{itemize}

\section*{When to Seek Medical Care}
See a doctor if you have several features of Cushing's syndrome, especially if you take steroids long-term or have a combination of weight gain, high blood pressure, bruising, and muscle weakness.

\textbf{Call 000 (or local emergency services) for:}
\begin{itemize}
    \item Severe infection or illness in a person taking long-term steroids, who may need extra steroid cover
    \item Sudden severe abdominal or back pain with collapse (possible adrenal crisis)
    \item Very high blood sugar with confusion or vomiting
    \item Chest pain or difficulty breathing
    \item A sudden severe headache or visual disturbance (possible pituitary tumour complications)
\end{itemize}

\section*{Diagnosis and Investigations}
Diagnosis requires confirming excess cortisol and then finding its source.

\begin{itemize}
    \item \textbf{History and examination:} features of cortisol excess and medication history
    \item \textbf{Salivary or urinary cortisol:} measured overnight or over 24 hours
    \item \textbf{Dexamethasone suppression test:} a low-dose steroid is given and cortisol measured
    \item \textbf{ACTH measurement:} distinguishes pituitary, adrenal, and ectopic causes
    \item \textbf{Imaging:} MRI of the pituitary and CT of the adrenal glands
    \item \textbf{Inferior petrosal sinus sampling:} a specialised test for difficult cases
    \item \textbf{Other tests:} blood sugar, blood pressure, bone density, and potassium
\end{itemize}

\section*{Management}
Treatment depends on the cause.

\begin{itemize}
    \item \textbf{Reducing steroid medicines:} gradual, supervised reduction of exogenous steroids, never abrupt withdrawal
    \item \textbf{Pituitary surgery:} transsphenoidal removal of the tumour, the first-line treatment for Cushing's disease
    \item \textbf{Adrenal surgery:} removal of an adrenal tumour
    \item \textbf{Radiotherapy:} for pituitary tumours not cured by surgery
    \item \textbf{Medicines:} drugs that block cortisol production or action
    \item \textbf{Removal of ectopic tumours:} surgery for the underlying cancer
    \item \textbf{Adrenalectomy:} removal of both adrenal glands in some cases
    \item \textbf{Management of complications:} blood pressure, diabetes, bone health, and infections
\end{itemize}

\section*{Complications}
\begin{itemize}
    \item Diabetes and its complications
    \item High blood pressure and cardiovascular disease
    \item Osteoporosis and fractures
    \item Muscle wasting and weakness
    \item Infections and poor wound healing
    \item Blood clots
    \item Depression and other mental health problems
    \item Growth failure in children
    \item Adrenal insufficiency after treatment, requiring lifelong replacement
    \item Complications of surgery and radiotherapy
\end{itemize}

\section*{Prognosis and Outcomes}
With successful treatment, most features of Cushing's syndrome gradually improve over months to years, and long-term outcomes are good. Weight normalises, blood pressure and blood sugar improve, and skin and muscle recover, although some changes may persist, especially if the condition has been longstanding. After treatment, people often develop temporary adrenal insufficiency and need replacement steroids while the adrenal glands recover. The outlook depends on the cause: pituitary and adrenal tumours have high cure rates, while ectopic Cushing's from cancer carries the prognosis of the underlying tumour.

\section*{Prevention and Self-Care}
\begin{itemize}
    \item Use steroid medicines at the lowest effective dose and for the shortest time possible
    \item Never stop long-term steroids suddenly; always taper under medical supervision
    \item Attend regular monitoring if on long-term steroids: blood pressure, blood sugar, and bone density
    \item Maintain a healthy diet and activity, with physiotherapy for muscle weakness
    \item Seek review for any new features of cortisol excess
    \item Carry a steroid card or warning bracelet if on long-term steroids
    \item Follow up with the endocrinology team after treatment
\end{itemize}

\section*{Sources and Further Reading}
The following reputable sources provide further information for healthcare students and patients seeking reliable, current advice about Cushing's syndrome.

\begin{itemize}
    \item Endocrine Society. \textit{Cushing's syndrome patient guide.} https://www.endocrine.org/
    \item Mayo Clinic. \textit{Cushing syndrome.} https://www.mayoclinic.org/diseases-conditions/cushing-syndrome/
    \item NHS. \textit{Cushing's syndrome.} https://www.nhs.uk/conditions/cushings-syndrome/
    \item MSD Manual. \textit{Cushing syndrome: professional reference.} https://www.msdmanuals.com/
    \item Healthdirect Australia. \textit{Cushing's syndrome.} https://www.healthdirect.gov.au/
    \item National Institute of Diabetes and Digestive and Kidney Diseases. \textit{Cushing's syndrome.} https://www.niddk.nih.gov/
\end{itemize}
